% One-page letter / statement template. Matches the CV letterhead.
\documentclass[11pt, letterpaper]{article}
\usepackage[T1]{fontenc}
\usepackage[utf8]{inputenc}
\usepackage{charter}
\usepackage[top=0.55in, bottom=0.45in, left=0.8in, right=0.8in]{geometry}
\usepackage[dvipsnames]{xcolor}
\definecolor{linkblue}{HTML}{1155CC}
\usepackage[colorlinks=true, urlcolor=linkblue, linkcolor=linkblue]{hyperref}
\usepackage{parskip}
\pagestyle{empty}
\setlength{\parindent}{0pt}
\linespread{1.0}

\newcommand{\letterhead}{%
  \begin{center}
    {\large\bfseries Rodrigo França Chaves}\\[2pt]
    {\small
      Chicago, IL \enspace$\cdot$\enspace
      \href{mailto:rchaves@uchicago.edu}{rchaves@uchicago.edu} \enspace$\cdot$\enspace
      +1 312 217 5526 \enspace$\cdot$\enspace
      \href{https://rodrigofch7.github.io}{rodrigofch7.github.io}
    }
  \end{center}
  \vspace{2pt}
  {\rule{\linewidth}{0.8pt}}
  \vspace{10pt}
}

\begin{document}
\letterhead

\textbf{The Brattle Group} \\
Energy Analyst (Economics), July 2027 \\
\textit{Preferred office: Washington, DC.}

\vspace{8pt}

Dear Hiring Committee,

I am applying for the Energy Analyst position in your Washington, DC office. I complete my Master's in Computational Analysis and Public Policy at the University of Chicago in June 2027, which lines up with a July start, and Washington is where I would most like to be: I spent this summer there with the World Bank and would be glad to return.

The work I would point to first is a study of how data center construction relates to what households nearby pay. I assembled the facility data myself, scraping 124 listed sites, confirming and geocoding 45, and dating each opening from air permit filings between 2000 and 2024, then matched those dates against ZIP-level home values and American Community Survey household electricity and water cost data across 660 metropolitan ZIP codes, and built the result into an interactive dashboard. The design is descriptive rather than causal and I was explicit about that throughout. It is the closest thing I have to your electricity practice: entry of large load onto a grid, and the question of who absorbs the cost.

My published work is in regulated utilities. I am first author of an article in \textit{Utilities Policy} asking whether transferring Brazil's water and wastewater systems to private operators improved public health, comparing propensity-matched municipalities before and after transfer between 1998 and 2021 using staggered difference-in-differences. It finds significant reductions in diarrheal morbidity among children under five, no effect on diseases unrelated to sanitation, and results that vary considerably across municipalities. I wrote the paper, assembled the data, and wrote the code. Ownership structure, service quality, and what a regulator can actually verify are the questions underneath it, and they are the questions regulatory proceedings turn on.

I am comfortable with the scale and the mess. At the World Bank I reconstructed four decades of Mumbai property tax and land-use policy into a single timeline and tested millions of assessment records for bunching at policy thresholds; I also built gridded panels from satellite rasters to measure construction responses to a new transit corridor, where I found that an inherited detection threshold was biasing estimates by neighborhood type. Separately I built a horizontal-equity audit of New York's property tax roll across 1.1 million properties. I work in R, Python, Stata, and Excel.

One thing I should state plainly: my undergraduate degree is in International Relations, not one of the quantitative disciplines you list. My quantitative training came afterwards and deliberately, through a postgraduate degree in Data Science at Insper and now the MSCAPP program, where I will also complete an energy-focused course this coming year, and I spent two of those years teaching econometrics and GIS to more than fifty students while working as a research assistant. I would rather you judge that on the work than on the label.

Thank you for your consideration.

\vspace{6pt}

Sincerely, \\
Rodrigo França Chaves

\end{document}
